\originalsection{第1次习题课}
\sourceinfo{MIT 18.04, Spring 2018, Recitation 1; 题面 \nolinkurl{mit18_04s18_recit1-handout.pdf}, PDF第1页; 解答 \nolinkurl{mit18_04s18_recit1-solutions.pdf}, PDF第1页. 课程作者 Jeremy Orloff, 习题课页内署名 Vishesh Jain, 2018, CC BY-NC-SA 4.0}
\begin{exercise}
\textbf{原题 Rec1.1。}\label{ass:rec1-1}
1.1. $\ee^{\log z}=z$ 是否成立? $\log(\ee^z)=z$ 是否成立? 1.2. 若知道 $\log z$ 的一个值, 其余值是什么?
\end{exercise}
\begin{solution}
据官方解答翻译补全; 1.1第二问原解答仅指向习题集2, 以下判断为编者补充 (AI). 设 $z=r\ee^{\ii\theta}\ne0$. 对任一个对数值 $\log r+\ii(\theta+2k\pi)$, 指数均为 $r\ee^{\ii\theta}=z$. 因此第一式对每个对数值成立. 若 $z=x+\ii y$, 则 $\log(\ee^z)$ 的全部值为 $x+\ii(y+2k\pi)=z+2k\pi\ii$, $k\in\Z$. 因而作为多值集合它包含 $z$, 却不等于单个 $z$; 若指定主支, 也未必取到 $z$, 例如 $z=2\pi\ii$ 时 $\Log(\ee^z)=0$. 采用 $-\pi<\Arg<\pi$ 的解析主支时, $-\pi<\Im z<\pi$ 才给出该等式; 边界须另按数值辐角约定处理. 1.2中, 若已知一个值 $L$, 所有值恰为 $L+2k\pi\ii$, $k\in\Z$. $z=0$ 没有复对数, 不属于这些陈述的定义域.
\end{solution}
\begin{exercise}
\textbf{原题 Rec1.2。}\label{ass:rec1-2}
设 $z_1=2\ee^{\ii\pi/3}$, $z_2=3+4\ii$. 2.1. 计算 $\log z_1$ 并给出主支值. 2.2. 对 $z_2$ 做同样计算.
\end{exercise}
\begin{solution}
据官方解答翻译补全. $\log z_1=\log2+\ii\pi/3+2k\pi\ii$, 主值为 $\log2+\ii\pi/3$. 令 $\theta=\arctan(4/3)\in(0,\pi/2)$, 则 $|z_2|=5$, $\log z_2=\log5+\ii\theta+2k\pi\ii$, 主值为 $\log5+\ii\theta$. 两式中的 $k$ 均遍历 $\Z$.
\end{solution}
\begin{exercise}
\textbf{原题 Rec1.3。}\label{ass:rec1-3}
3.1. $z^a$ 是单值还是多值? 为什么? 3.2. 假设 $z\ne0$, $a$ 为整数时呢? 3.3. $a$ 为实数时, $z$ 的所有 $a$ 次幂有什么共同性质? 3.4. $a$ 为纯虚数时呢?
\end{exercise}
\begin{solution}
据官方解答翻译补全, 3.1的例外说明为编者补充 (AI). 在 $z\ne0$ 下固定 $z=r\ee^{\ii\theta}$. 其幂值为 $\ee^{a(\log r+\ii\theta+2k\pi\ii)}$, 相邻值的比为 $\ee^{2\pi\ii a}$. 因而一般为多值; 当且仅当 $a\in\Z$ 时全部值相同. 这说明原解答“多值”应理解为一般情形, 不能排除整数例外. 3.2由 $\ee^{2k\pi\ii a}=1$ 得到单值, 与整数乘方及负整数倒数一致. 3.3若 $a\in\R$, 每个值的模均为 $r^a$. 3.4若 $a=\ii b$, $b\in\R$, 则各值为 $\ee^{-b(\theta+2k\pi)}\ee^{\ii b\log r}$, 它们位于同一条从原点出发的射线上, 因而主辐角相同. 当 $b=0$ 时各值均为1.
\end{solution}
\begin{exercise}
\textbf{原题 Rec1.4。}\label{ass:rec1-4}
设 $z_1=2\ee^{\ii\pi/3}$, $z_2=3+4\ii$. 4.1. 计算 $z_1^{z_2}$. 4.2. 计算 $z_1^{1/4}$, 有多少个不同值? 在复平面标出全部值.
\end{exercise}
\begin{solution}
据官方解答翻译补全; 图为编者补充 (AI). 4.1按多值幂的定义,
\[
z_1^{z_2}=\ee^{(3+4\ii)(\log2+\ii\pi/3+2k\pi\ii)}
=\ee^{3\log2-4\pi/3-8k\pi}\ee^{\ii(4\log2+\pi+6k\pi)},\quad k\in\Z.
\]
最后的 $\ee^{6k\pi\ii}=1$, 但模随 $k$ 改变, 所以有无穷多个值. 主值对应 $k=0$. 4.2四次根为 $2^{1/4}\ee^{\ii(\pi/12+k\pi/2)}$, $k=0,1,2,3$. 指标相差4时重复, 这四个值互异, 是半径 $2^{1/4}$ 的圆上正方形的顶点.
\begin{center}
\begin{tikzpicture}[scale=1.3,>=Stealth]
\draw[->] (-1.6,0)--(1.6,0) node[right]{$\Re w$};\draw[->] (0,-1.6)--(0,1.6) node[above]{$\Im w$};
\draw[dashed] (0,0) circle (1.1892);
\foreach \k in {0,1,2,3}{\pgfmathsetmacro{\a}{15+90*\k}\fill ({1.1892*cos(\a)},{1.1892*sin(\a)}) circle(1.3pt) node[shift={(\a:0.35)}]{$w_{\k}$};}
\draw (0.6,0) arc(0:15:0.6);\node at(.75,.15){$\pi/12$};
\end{tikzpicture}
\end{center}
\end{solution}
